% Generated draft: requires mathematical and source review.
\documentclass[11pt]{article}
\usepackage[margin=25mm]{geometry}
\usepackage{fontspec}
\setmainfont{Latin Modern Roman}
\usepackage{amsmath,amssymb,mathtools,unicode-math}
\setmathfont{Latin Modern Math}
\usepackage{enumitem,needspace,xurl,hyperref}
\setlength{\parindent}{0pt}
\setlength{\parskip}{5pt}
\setlength{\emergencystretch}{3em}
\newcommand{\N}{\mathbb{N}}
\newcommand{\Z}{\mathbb{Z}}
\newcommand{\Q}{\mathbb{Q}}
\newcommand{\R}{\mathbb{R}}
\newcommand{\C}{\mathbb{C}}
\newcommand{\F}{\mathbb{F}}
\newcommand{\A}{\mathbb{A}}
\newcommand{\PP}{\mathbb{P}}
\newcommand{\E}{\mathbb{E}}
\newcommand{\eps}{\varepsilon}
\newcommand{\dd}{\,\mathrm{d}}
\begin{document}
\section*{UnsolvedMath ID: 701129}
\textbf{Source problem code:} catalogue entry
\subsection{Definitions and mathematical statement}
For a positive integer \(m\) and a prime \(p\), let \(R(m,p)\) denote the largest finite group generated by \(m\) elements and having exponent \(p\), as in the supplied source. Interpret ``largest'' in the usual universal sense: every finite group generated by at most \(m\) elements and satisfying \(g^p=1\) for every element \(g\) is a quotient of \(R(m,p)\). The source does not explicitly specify this universal property.

For a group \(G\), define its lower central series by \(\gamma_1(G)=G\) and \(\gamma_{i+1}(G)=[\gamma_i(G),G]\), where \([H,G]\) is generated by the commutators \(h^{-1}g^{-1}hg\) with \(h\in H\) and \(g\in G\). For a nontrivial nilpotent group \(G\), its nilpotency class \(\operatorname{cl}(G)\) is the least positive integer \(c\) such that \(\gamma_{c+1}(G)=1\). Write \(c_p(m)=\operatorname{cl}(R(m,p))\).

The following questions use the interpretation that the prime \(p\) is fixed while \(m\) varies, so the bounding constants and polynomial degree may depend on \(p\). The source does not explicitly settle whether bounds uniform in \(p\) are intended.
\begin{enumerate}
\item[(a)] Is it true that for every prime \(p\) there exist a real constant \(C_p>0\) and an integer \(d_p\geq 0\) such that
\[
\forall m\in\Z_{\geq 1},\qquad c_p(m)\leq C_p m^{d_p}?
\]
The supplied source states that a polynomial bound holds for \(p=2,3,5,7\).
\item[(b)] Is it true that for every prime \(p\) there exist real constants \(A_p,B_p\geq 0\) such that
\[
\forall m\in\Z_{\geq 1},\qquad c_p(m)\leq A_p m+B_p?
\]
The supplied source states that a linear bound holds for \(p=2,3,5\).
\item[(c)] In particular, do there exist real constants \(A,B\geq 0\) such that
\[
\forall m\in\Z_{\geq 1},\qquad \operatorname{cl}(R(m,7))\leq Am+B?
\]
\end{enumerate}
M. R. Vaughan-Lee expresses the guess that the answers to (a) and (b) are negative for general primes, whereas the answer to (c) is positive. These are attributed guesses, not conclusions established in this draft.

As background, the source also uses \(R(m,p^k)\) for the analogous largest finite \(m\)-generator group of exponent \(p^k\). It attributes to Mike Newman an argument giving the following order bound. For every prime \(p\), integers \(m\geq 2\) and \(k\geq 2\), with the additional restriction \(k\geq 3\) when \(p=2\), define
\[
T_0=m,\qquad T_{j+1}=p^{T_j}\quad(0\leq j<k).
\]
Then the supplied source asserts
\[
|R(m,p^k)|\geq T_k,
\]
so the tower contains exactly \(k\) occurrences of \(p\), with \(m\) as the top exponent. This is supplied background concerning group order, distinct from the questions about nilpotency class.
\subsection{Short English statement}
For each fixed prime \(p\), ask whether the nilpotency class of the largest finite \(m\)-generator group of exponent \(p\) has a polynomial bound, or even a linear bound, as \(m\) varies. The specific remaining question highlighted by the source is a linear bound for exponent \(7\). Dependence of the bounds on \(p\) is an interpretation requiring review. The registry labels the problem open; this label and the stated special cases are supplied source information, not independently verified assessments.
\subsection{Sources}
\begin{itemize}
\item Supplied local UnsolvedMath registry record: permanent ID \(701129\), exact source problem code \texttt{catalogue entry}, title ``Kourovka Notebook Problem 15.41.'' The record labels the problem \texttt{open}; its supplied creation and update dates are 2026-10-16. Upstream revision: \texttt{d371ee917ac9647fbb06cd6fc529deec6f3dda82}.
\item \emph{The Kourovka Notebook (Unsolved Problems in Group Theory)}, Problem 15.41, proposed by M. R. Vaughan-Lee in the 15th issue (2002). The supplied record identifies the 21st issue edition (2026), arXiv identifier \texttt{1401.0300v48}. The collection description names E. I. Khukhro and V. D. Mazurov as editors. These bibliographic identifications are supplied attribution and have not been independently checked.
\item M. Vaughan-Lee and E. I. Zelmanov, \emph{J. Austral. Math. Soc. (A)} \textbf{67}, no. 2 (1999), 261--271. The supplied problem cites this reference for the prime-power exponent order bound and credits its argument to Mike Newman. No article title or external link is supplied.
\end{itemize}
\end{document}
